\ifdefined\gkver\else\def\gkver{student}\fi
\documentclass[\gkver]{gaokaozhenti}
\newcommand{\ccnum}[1]{{\fontspec{Noto Sans CJK JP}#1}}
\begin{document}

\chapter{2009年宁夏理综}
%% paperType: 理综物理部分

\begin{enumerate}
\item
%% number: 14
%% paperName: 2009年宁夏理综第 14 题 6 分
%% typeId: 2
%% type: 多选题
%% chapter: 其他
%% point: 物理学史
%% method: 无
%% score: 6
%% degree: 650
%% duplicateId: 0
%% body:
在力学理论建立的过程中，有许多伟大的科学家做出了贡献。关于科学家和他们的贡献，下列说法正确的是\xzanswer{BD}

\fourchoices[answer=BD]
{伽利略发现了行星运动的规律}
{卡文迪许通过实验测出了引力常量}
{牛顿最早指出力不是维持物体运动的原因}
{笛卡尔对牛顿第一定律的建立做出了贡献}

%% answer: BD
%% memo:
\memoanswer{
行星运动定律是开普勒发现的，A 错误；B 正确；伽利略最早指出力不是维持物体运动的原因，C 错误；D 正确。
}

\item
%% number: 15
%% paperName: 2009年宁夏理综第 15 题 6 分
%% typeId: 1
%% type: 单选题
%% chapter: 曲线运动
%% point: 天体运动
%% method: 无
%% score: 6
%% degree: 650
%% duplicateId: 0
%% body:
地球和木星绕太阳运行的轨道都可以看作是圆形的。已知木星的轨道半径约为地球轨道半径的 5.2 倍，则木星与地球绕太阳运行的线速度之比约为\xzanswer{B}

\fourchoices[answer=B]
{0.19}
{0.44}
{2.3}
{5.2}

%% answer: B
%% memo:
\memoanswer{
天体的运动满足万有引力充当向心力即 $G\frac{Mm}{R^{2}}=m\frac{v^{2}}{R}$，可知 $v=\sqrt{\frac{GM}{R}}$，可见木星与地球绕太阳运行的线速度之比 $\frac{v_{\text{木}}}{v_{\text{地}}}=\sqrt{\frac{R_{\text{地}}}{R_{\text{木}}}}=\sqrt{\frac{1}{5.2}}\approx0.44$，B 正确。
}

\item
%% number: 16
%% paperName: 2009年宁夏理综第 16 题 6 分
%% typeId: 1
%% type: 单选题
%% chapter: 电磁感应
%% point: 切割磁感线电动势
%% method: 建立模型
%% score: 6
%% degree: 450
%% duplicateId: 0
%% body:
医生做某些特殊手术时，利用电磁血流计来监测通过动脉的血流速度。电磁血流计由一对电极 a 和 b 以及磁极 N 和 S 构成，磁极间的磁场是均匀的。使用时，两电极 a、b 均与血管壁接触，两触点的连线、磁场方向和血流速度方向两两垂直，如图所示。由于血液中的正负离子随血流一起在磁场中运动，电极 a、b 之间会有微小电势差。在达到平衡时，血管内部的电场可看作是匀强电场，血液中的离子所受的电场力和磁场力的合力为零。在某次监测中，两触点的距离为 $3.0\Umm$，血管壁的厚度可忽略，两触点间的电势差为 $160\UuV$，磁感应强度的大小为 $0.040\UT$。则血流速度的近似值和电极 a、b 的正负为\xzanswer{A}

\onepicture[w=5.5cm, align=right]{figs/16.png}

\fourchoices[answer=A]
{$1.3\Ums$，a 正、b 负}
{$2.7\Ums$，a 正、b 负}
{$1.3\Ums$，a 负、b 正}
{$2.7\Ums$，a 负、b 正}

%% answer: A
%% memo:
\memoanswer{
依据左手定则，正离子在磁场中受到洛伦兹力作用向上偏，负离子在磁场中受到洛伦兹力作用向下偏，因此电极 a、b 的正负为 a 正、b 负；当稳定时，血液中的离子所受的电场力和磁场力的合力为零，则 $qE=qvB$，可得 $v=\frac{E}{B}=\frac{U}{Bd}=\frac{160\times10^{-6}}{0.04\times3\times10^{-3}}\Ums\approx1.3\Ums$，A 正确。
}

\item
%% number: 17
%% paperName: 2009年宁夏理综第 17 题 6 分
%% typeId: 2
%% type: 多选题
%% chapter: 功和能
%% point: 功率
%% method: 无
%% score: 6
%% degree: 450
%% duplicateId: 0
%% body:
质量为 $m$ 的物体静止在光滑水平面上，从 $t=0$ 时刻开始受到水平力的作用。力的大小 $F$ 与时间 $t$ 的关系如图所示，力的方向保持不变，则\xzanswer{BD}

\onepicture[w=5.0cm, align=right]{figs/17.png}

\fourchoices[answer=BD]
{$3t_{0}$ 时刻的瞬时功率为 $\frac{5F_0^{2}t_0}{m}$}
{$3t_{0}$ 时刻的瞬时功率为 $\frac{15F_0^{2}t_0}{m}$}
{在 $t=0$ 到 $3t_{0}$ 这段时间内，水平力的平均功率为 $\frac{23F_0^{2}t_0}{4m}$}
{在 $t=0$ 到 $3t_{0}$ 这段时间内，水平力的平均功率为 $\frac{25F_0^{2}t_0}{6m}$}

%% answer: BD
%% memo:
\memoanswer{
$0\sim2t_{0}$ 内物体的加速度为 $a=\frac{F_0}{m}$，$2t_{0}$ 时刻的速度为 $v_{1}=a\cdot2t_{0}=\frac{2F_0t_0}{m}$，在 $3t_{0}$ 时刻的瞬时速度 $v_{2}=v_{1}+3a\cdot t_{0}=\frac{5F_0t_0}{m}$，则 $3t_{0}$ 时刻的瞬时功率为 $P=3F_0\cdot\frac{5F_0t_0}{m}=\frac{15F_0^{2}t_0}{m}$，A 错误，B 正确；在 $t=0$ 到 $3t_{0}$ 这段时间内，由动能定理可得 $W=\Delta E_k=\frac{1}{2}mv_{2}^{2}=\frac{25F_0^{2}t_0^{2}}{2m}$，则这段时间内的平均功率 $\overline{P}=\frac{25F_0^{2}t_0}{6m}$，D 正确。
}

\item
%% number: 18
%% paperName: 2009年宁夏理综第 18 题 6 分
%% typeId: 1
%% type: 单选题
%% chapter: 电场
%% point: 电势差与电场强度的关系
%% method: 无
%% score: 6
%% degree: 450
%% duplicateId: 0
%% body:
空间有一均匀强电场，在电场中建立如图所示的直角坐标系 O-$xyz$，M、N、P 为电场中的三个点，M 点的坐标（0，$a$，0），N 点的坐标为（$a$，0，0），P 点的坐标为（$a$，$\frac{a}{2}$，$\frac{a}{2}$）。已知电场方向平行于直线 MN，M 点电势为 0，N 点电势为 $1\UV$，则 P 点的电势为\xzanswer{D}

\onepicture[w=5.0cm, align=right]{figs/18.png}

\fourchoices[answer=D]
{$\frac{\sqrt{2}}{2}\UV$}
{$\frac{\sqrt{3}}{2}\UV$}
{$\frac{1}{4}\UV$}
{$\frac{3}{4}\UV$}

%% answer: D
%% memo:
\memoanswer{
将立体图画成平面图，如图所示，可见 P 点沿电场线方向为 MN 的四等分线，故 P 点的电势为 $\frac{3}{4}\UV$，D 正确。

\onepicture[w=5.0cm]{figs/18_an.png}
}

\item
%% number: 19
%% paperName: 2009年宁夏理综第 19 题 6 分
%% typeId: 1
%% type: 单选题
%% chapter: 电磁感应
%% point: 切割磁感线电动势
%% method: 无
%% score: 6
%% degree: 450
%% duplicateId: 0
%% body:
如图所示，一导体圆环位于纸面内，O 为圆心。环内两个圆心角为 $90^{\circ}$ 的扇形区域内分别有匀强磁场，两磁场磁感应强度的大小相等，方向相反且均与纸面垂直。导体杆 OM 可绕 O 转动，M 端通过滑动触点与圆环良好接触。在圆心和圆环间连有电阻 $R$。杆 OM 以匀角速度 $\omega$ 逆时针转动，$t=0$ 时恰好在图示位置。规定从 a 到 b 流经电阻 $R$ 的电流方向为正，圆环和导体杆的电阻忽略不计，则杆从 $t=0$ 开始转动一周的过程中，电流随 $\omega t$ 变化的图象是\xzanswer{C}

\onepicture[w=5.0cm]{figs/19a.png}

\fourchoices[answer=C, ispicture=true, h=2.5cm]
{figs/19b.png}
{figs/19c.png}
{figs/19d.png}
{figs/19e.png}

%% answer: C
%% memo:
\memoanswer{
依据右手定则，可知在 $0\sim\frac{\pi}{2}$ 内，电流方向 M 到 O，在电阻 $R$ 内则是由 b 到 a，为负值，且大小为 $I=\frac{\frac{1}{2}BL\omega^{2}}{R}$ 为一定值；$\frac{\pi}{2}\sim\pi$ 内没有感应电流；$\pi\sim\frac{3}{2}\pi$ 内电流的方向相反，即沿正方向；$\frac{3\pi}{2}\sim2\pi$ 内没有感应电流，因此 C 正确。
}

\item
%% number: 20
%% paperName: 2009年宁夏理综第 20 题 6 分
%% typeId: 2
%% type: 多选题
%% chapter: 运动和力
%% point: 传送带
%% method: 无
%% score: 6
%% degree: 450
%% duplicateId: 0
%% body:
如图所示，一足够长的木板静止在光滑水平面上，一物块静止在木板上，木板和物块间有摩擦。现用水平力向右拉木板，当物块相对木板滑动了一段距离但仍有相对运动时，撤掉拉力，此后木板和物块相对于水平面的运动情况为\xzanswer{BC}

\onepicture[w=5.5cm, align=right]{figs/20.png}

\fourchoices[answer=BC]
{物块先向左运动，再向右运动}
{物块向右运动，速度逐渐增大，直到做匀速运动}
{木板向右运动，速度逐渐变小，直到做匀速运动}
{木板和物块的速度都逐渐变小，直到为零}

%% answer: BC
%% memo:
\memoanswer{
对于物块，由于运动过程中与木板存在相对滑动，且始终相对木板向左运动，因此木板对物块的摩擦力向右，所以物块相对地面向右运动，且速度不断增大，直至相对静止而做匀速直线运动，B 正确；对于木板，由作用力与反作用力可知受到物块给它的向左的摩擦力作用，则木板的速度不断减小，直到二者相对静止，而做匀速直线运动，C 正确；由于水平面光滑，所以不会停止，D 错误。
}

\item
%% number: 21
%% paperName: 2009年宁夏理综第 21 题 6 分
%% typeId: 2
%% type: 多选题
%% chapter: 运动和力
%% point: 牛顿第二定律
%% method: 极值
%% score: 6
%% degree: 250
%% duplicateId: 0
%% body:
水平地面上有一木箱，木箱与地面之间的动摩擦因数为 $\mu$（$0<\mu<1$）。现对木箱施加一拉力 $F$，使木箱做匀速直线运动。设 $F$ 的方向与水平面夹角为 $\theta$，如图，在 $\theta$ 从 0 逐渐增大到 $90^{\circ}$ 的过程中，木箱的速度保持不变，则\xzanswer{AC}

\onepicture[w=4.5cm, align=right]{figs/21.png}

\fourchoices[answer=AC]
{$F$ 先减小后增大}
{$F$ 一直增大}
{$F$ 的功率减小}
{$F$ 的功率不变}

%% answer: AC
%% memo:
\memoanswer{
由于木箱的速度保持不变，因此木箱始终处于平衡状态，受力分析如图所示，则由平衡条件得 $mg=N+F\sin\theta$，$f=\mu N=F\cos\theta$，两式联立解得 $F=\frac{\mu mg}{\cos\theta+\mu\sin\theta}=\frac{\mu mg}{\sqrt{1+\mu^{2}}\sin(\theta+\alpha)}$，可见 $F$ 有最小值，所以 $F$ 先减小后增大，A 正确，B 错误；$F$ 的功率 $P=Fv\cos\theta=\frac{\mu mgv\cos\theta}{\cos\theta+\mu\sin\theta}=\frac{\mu mgv}{1+\mu\tan\theta}$，可见在 $\theta$ 从 0 逐渐增大到 $90^{\circ}$ 的过程中 $\tan\theta$ 逐渐增大，则功率 $P$ 逐渐减小，C 正确，D 错误。
}

\item
%% number: 22
%% paperName: 2009年宁夏理综第 22 题 4 分
%% typeId: 4
%% type: 实验题
%% chapter: 直线运动
%% point: 游标卡尺和螺旋测微仪
%% method: 无
%% score: 4
%% degree: 650
%% duplicateId: 0
%% body:
某同学用游标卡尺测量一圆柱体的长度 $l$，用螺旋测微器测量该圆柱体的直径 $d$，示数如图。由图可读出 $l=$ \tkanswer[2cm]{} $\Ucm$，$d=$ \tkanswer[2cm]{} $\Umm$。

\onepicture[w=11cm, align=right]{figs/22.png}

%% answer:
\jdanswer{
$2.25\Ucm$，$6.860\Umm$。
}
%% memo:
\memoanswer{
游标卡尺的读数 $x=22\Umm+0.1\Umm\times5=22.5\Umm=2.25\Ucm$；螺旋测微器的读数 $y=6.5\Umm+0.01\Umm\times36.0=6.860\Umm$。
}

\item
%% number: 23
%% paperName: 2009年宁夏理综第 23 题 11 分
%% typeId: 4
%% type: 实验题
%% chapter: 电路
%% point: 电容
%% method: 无
%% score: 11
%% degree: 450
%% duplicateId: 0
%% body:
青岛奥运会帆船赛场采用风力发电给蓄电池充电，为路灯提供电能。用光敏电阻作为传感器控制路灯电路的开关，实现自动控制。光敏电阻的阻值随照射光的强弱而变化，作为简化模型，可以近似认为，照射光较强（如白天）时电阻几乎为 0；照射光较弱（如黑天）时电阻接近于无穷大。利用光敏电阻作为传感器，借助电磁开关，可以实现路灯自动在白天关闭，黑天打开。电磁开关的内部结构如图所示。1、2 两接线柱之间是励磁线圈，3、4 两接线柱分别与弹簧片和触点连接。当励磁线圈中电流大于 $50\UmA$ 时，电磁铁吸合铁片，弹簧片和触点分离，3、4 断开；电流小于 $50\UmA$ 时，3、4 接通。励磁线圈中允许通过的最大电流为 $100\UmA$。

\onepicture[w=4.5cm, align=right]{figs/23a.png}

\begin{enumerate}
	\item
	利用以下器材设计一个自动控制路灯的电路，画出电路原理图。

	光敏电阻 $R_{1}$，符号；灯泡 L，额定功率 $40\UW$，额定电压 $36\UV$，符号；保护电阻 $R_{2}$，符号；电磁开关，符号；蓄电池 E，电压 $36\UV$，内阻很小；开关 S，导线若干。

	\fourpicture[showlabel=false, w=1.2cm]
	{figs/23b.png}
	{figs/23c.png}
	{figs/23d.png}
	{figs/23e.png}

	\item
	回答下列问题：

	\begin{enumerate}
		\item
		如果励磁线圈的电阻为 $200\UO$，励磁线圈允许加的最大电压为 \tkanswer[1.2cm]{} $\UV$，保护电阻 $R_{2}$ 的阻值范围为 \tkanswer[2cm]{} $\UO$。

		\item
		在有些应用电磁开关的场合，为了安全，往往需要在电磁铁吸合铁片时，接线柱 3、4 之间从断开变为接通。为此，电磁开关内部结构应如何改造？请结合本题中电磁开关内部结构图说明。

		答：\tkanswer[6cm]{}。

		\item
		任意举出一个其它的电磁铁应用的例子。

		答：\tkanswer[4cm]{}。
	\end{enumerate}
\end{enumerate}

%% answer:
\jdanswer{
\begin{enumerate}
	\item
	电路原理图见解析。

	\item
	\begin{enumerate}
		\item
		$20\UV$；$160\sim520\UO$

		\item
		把触点从弹簧片右侧移到弹簧片左侧，保证当电磁铁吸合铁片时，3、4 之间接通；不吸合时，3、4 之间断开。

		\item
		电磁起重机
	\end{enumerate}
\end{enumerate}
}
%% memo:
\memoanswer{
（1）要使光敏电阻能够对电路进行控制，且有光照时路灯熄灭，光敏电阻应与 1、2 串联，3、4 与路灯串联；则电路图如图所示。

\onepicture[w=6.0cm]{figs/23_an.png}

（2）①由 $U=IR$ 得励磁线圈允许加的最大电压为 $U=I_{m}R=0.1\times200\UV=20\UV$；依据允许通过励磁线圈的电流最大值和最小值计算得 $R_{2\min}=\frac{16}{0.1}\UO=160\UO$，$R_{2\max}=\frac{16}{0.05}\UO=320\UO$，因此保护电阻 $R_{2}$ 的阻值范围为 $160\sim320\UO$；②把触点从弹簧片右侧移到弹簧片左侧，保证当电磁铁吸合铁片时，3、4 之间接通；不吸合时，3、4 之间断开。

③电磁起重机
}

\item
%% number: 24
%% paperName: 2009年宁夏理综第 24 题 14 分
%% typeId: 6
%% type: 计算题
%% chapter: 运动和力
%% point: 牛顿定律的应用
%% method: 无
%% score: 14
%% degree: 450
%% duplicateId: 0
%% body:
冰壶比赛是在水平冰面上进行的体育项目，比赛场地示意如图。比赛时，运动员从起滑架处推着冰壶出发，在投掷线 AB 处放手让冰壶以一定的速度滑出，使冰壶的停止位置尽量靠近圆心 O。为使冰壶滑行得更远，运动员可以用毛刷擦冰壶运行前方的冰面，使冰壶与冰面间的动摩擦因数减小。设冰壶与冰面间的动摩擦因数为 $\mu_{1}=0.008$，用毛刷擦冰面后动摩擦因数减少至 $\mu_{2}=0.004$。在某次比赛中，运动员使冰壶 C 在投掷线中点处以 $2\Ums$ 的速度沿虚线滑出。为使冰壶 C 能够沿虚线恰好到达圆心 O 点，运动员用毛刷擦冰面的长度应为多少？（$g$ 取 $10\Umsq$）

\onepicture[w=10cm, align=right]{\includesvg{figs/24}}

%% answer:
\jdanswer{
$10\Um$
}
%% memo:
\memoanswer{
设冰壶在未被毛刷擦过的冰面上滑行的距离为 $S_{1}$，所受摩擦力的大小为 $f_{1}$；在被毛刷擦过的冰面上滑行的距离为 $S_{2}$，所受摩擦力的大小为 $f_{2}$。则有 $S_{1}+S_{2}=S$ ①，式中 $S$ 为投掷线到圆心 O 的距离。

$f_{1}=\mu_{1}mg$ ②，$f_{2}=\mu_{2}mg$ ③。

设冰壶的初速度为 $v_{0}$，由功能关系，得 $f_{1}\cdot S_{1}+f_{2}\cdot S_{2}=\frac{1}{2}mv_{0}^{2}$ ④。

联立以上各式，解得 $S_{2}=\frac{2\mu_{1}gS-v_{0}^{2}}{2g(\mu_{1}-\mu_{2})}$ ⑤。

代入数据得 $S_{2}=10\Um$ ⑥。
}

\item
%% number: 25
%% paperName: 2009年宁夏理综第 25 题 18 分
%% typeId: 6
%% type: 计算题
%% chapter: 磁场
%% point: 带电粒子在组合场中的运动
%% method: 无
%% score: 18
%% degree: 450
%% duplicateId: 0
%% body:
如图所示，在第一象限有一均强电场，场强大小为 $E$，方向与 $y$ 轴平行；在 $x$ 轴下方有一均强磁场，磁场方向与纸面垂直。一质量为 $m$、电荷量为 $-q$（$q>0$）的粒子以平行于 $x$ 轴的速度从 $y$ 轴上的 P 点处射入电场，在 $x$ 轴上的 Q 点处进入磁场，并从坐标原点 O 离开磁场。粒子在磁场中的运动轨迹与 $y$ 轴交于 M 点。已知 OP$=l$，OQ$=2\sqrt{3}l$。不计重力。求：

\begin{enumerate}
	\item
	M 点与坐标原点 O 间的距离；
	\item
	粒子从 P 点运动到 M 点所用的时间。
\end{enumerate}

\onepicture[w=5.0cm, align=right]{figs/25.png}

%% answer:
\jdanswer{
\begin{enumerate}
	\item
	$MO=6l$
	\item
	$t=\left(\frac{\sqrt{3}}{2}\pi+1\right)\sqrt{\frac{2ml}{qE}}$
\end{enumerate}
}
%% memo:
\memoanswer{
（1）带电粒子在电场中做类平抛运动，在 $y$ 轴负方向上做初速度为零的匀加速运动，设加速度的大小为 $a$；在 $x$ 轴正方向上做匀速直线运动，设速度为 $v_{0}$，粒子从 P 点运动到 Q 点所用的时间为 $t_{1}$，进入磁场时速度方向与 $x$ 轴正方向的夹角为 $\theta$，则 $a=\frac{qE}{m}$ ①，$t_{1}=\sqrt{\frac{2y_{0}}{a}}$ ②，$v_{0}=\frac{x_{0}}{t_{1}}$ ③，其中 $x_{0}=2\sqrt{3}l$，$y_{0}=l$。又有 $\tan\theta=\frac{at_{1}}{v_{0}}$ ④。联立②③④式，得 $\theta=30^{\circ}$ ⑤。因为 M、O、Q 点在圆周上，$\angle MOQ=90^{\circ}$，所以 MQ 为直径。从图中的几何关系可知 $R=2\sqrt{3}l$ ⑥，$MO=6l$ ⑦。

\onepicture[w=5.5cm]{figs/25_an.png}

（2）设粒子在磁场中运动的速度为 $v$，从 Q 到 M 点运动的时间为 $t_{2}$，则有 $v=\frac{v_{0}}{\cos\theta}$ ⑧，$t_{2}=\frac{\pi R}{v}$ ⑨。带电粒子自 P 点出发到 M 点所用的时间为 $t=t_{1}+t_{2}$ ⑩。联立①②③⑤⑥⑧⑨⑩式，并代入数据得 $t=\left(\frac{\sqrt{3}}{2}\pi+1\right)\sqrt{\frac{2ml}{qE}}$\ccnum{⑪}。
}

\item
%% number: 33
%% paperName: 2009年宁夏理综第 33 题 10 分
%% typeId: 6
%% type: 计算题
%% chapter: 功和能
%% point: 功
%% method: 无
%% score: 10
%% degree: 650
%% duplicateId: 0
%% body:
\textbf{（1）}（5 分）常见的传动方式有 \tkanswer[1.6cm]{}、\tkanswer[1.6cm]{}、\tkanswer[1.6cm]{} 和齿轮传动等。齿轮传动的传动比是主动轮与 \tkanswer[1.6cm]{} 的转速之比，传动比等于 \tkanswer[1.4cm]{} 与 \tkanswer[1.4cm]{} 的齿数之比。

\textbf{（2）}（10 分）液压千斤顶是利用密闭容器内的液体能够把液体所受到的压强行各个方向传递的原理制成的。图为一小型千斤顶的结构示意图。大活塞的直径 $D_{1}=20\Ucm$，小活塞 B 的直径 $D_{2}=5\Ucm$，手柄的长度 OC$=50\Ucm$，小活塞与手柄的连接点到转轴 O 的距离 OD$=10\Ucm$。现用此千斤顶使质量 $m=4\times10^{3}\Ukg$ 的重物升高了 $h=10\Ucm$。$g$ 取 $10\Umsq$，求：

\textbf{（i）}若此千斤顶的效率为 $80\%$，在这一过程中人做的功为多少？

\textbf{（ii）}若此千斤顶的效率为 $100\%$，当重物上升时，人对手柄的作用力 $F$ 至少要多大？

\onepicture[w=7.0cm, align=right]{figs/33.png}

%% answer:
\jdanswer{
\textbf{（1）}链传动，带传动，液压传动，从动轮，从动轮，主动轮。

\textbf{（2）}（i）$W_{2}=5\times10^{3}\UJ$；（ii）$F=500\UN$。
}
%% memo:
\memoanswer{
（1）常见的传动方式有链传动、带传动、液压传动和齿轮传动等；齿轮传动的传动比是主动轮与从动轮的转速之比，传动比等于从动轮与主动轮的齿数之比。

（2）（i）将重物托起 $h$ 需要做的功 $W_{1}=mgh$ ①。设人对手柄做的功为 $W_{2}$，则千斤顶的效率为 $\eta=\frac{W_{1}}{W_{2}}$ ②。代入数据可得 $W_{2}=5\times10^{3}\UJ$ ③。

（ii）设大活塞的面积为 $S_{1}$，小活塞的面积为 $S_{2}$，作用在小活塞上的压力为 $F_{1}$，当千斤顶的效率为 $100\%$ 时，有 $\frac{mg}{F_{1}}=\frac{S_{1}}{S_{2}}$ ④，$\frac{S_{1}}{S_{2}}=\frac{D_{1}^{2}}{D_{2}^{2}}$ ⑤。当 $F_{1}$ 和 $F$ 都与杠杆垂直时，手对杠杆的压力最小。利用杠杆原理，有 $F_{1}\cdot OD=F\cdot OC$ ⑥。由④⑤⑥式得 $F=500\UN$ ⑦。
}

\item
%% number: 34
%% paperName: 2009年宁夏理综第 34 题 10 分
%% typeId: 6
%% type: 计算题
%% chapter: 气体性质
%% point: 气体多过程
%% method: 无
%% score: 10
%% degree: 450
%% duplicateId: 0
%% body:
\textbf{（1）}（5 分）带有活塞的汽缸内封闭一定量的理想气体。气体开始处于状态 a，然后经过过程 ab 到达状态 b 或经过过程 ac 到状态 c，b、c 状态温度相同，如 $V-T$ 图所示。设气体在状态 b 和状态 c 的压强分别为 $p_{b}$ 和 $p_{c}$，在过程 ab 和 ac 中吸收的热量分别为 $Q_{ab}$ 和 $Q_{ac}$，则\xzanswer{C}（填入选项前的字母，有填错的不得分）。

\onepicture[w=4.5cm]{figs/34a.png}

\fourchoices[answer=C]
{$p_{b}>p_{c}$，$Q_{ab}>Q_{ac}$}
{$p_{b}>p_{c}$，$Q_{ab}<Q_{ac}$}
{$p_{b}<p_{c}$，$Q_{ab}>Q_{ac}$}
{$p_{b}<p_{c}$，$Q_{ab}<Q_{ac}$}

\textbf{（2）}（10 分）图中系统由左右两个侧壁绝热、底部、截面均为 $S$ 的容器组成。左容器足够高，上端敞开，右容器上端由导热材料封闭。两个容器的下端由可忽略容积的细管连通。容器内两个绝热的活塞 A、B 下方封有氮气，B 上方封有氢气。大气的压强 $p_{0}$，温度为 $T_{0}=273\UK$，两个活塞因自身重量对下方气体产生的附加压强均为 $0.1p_{0}$。系统平衡时，各气体柱的高度如图所示。现将系统的底部浸入恒温热水槽中，再次平衡时 A 上升了一定的高度。用外力将 A 缓慢推回第一次平衡时的位置并固定，第三次达到平衡后，氢气柱高度为 $0.8h$。氮气和氢气均可视为理想气体。求：

\begin{enumerate}
	\item
	第二次平衡时氮气的体积；
	\item
	水的温度。
\end{enumerate}

\onepicture[align=right]{figs/34.png}

%% answer:
\jdanswer{
\textbf{（1）}C

\textbf{（2）}
\begin{enumerate}
	\item
	$2.7hS$
	\item
	$368.55\UK$
\end{enumerate}
}
%% memo:
\memoanswer{
（1）由 $V-T$ 图可知，过程 ab 为等压过程，$p_{b}=p_{a}$；过程 ac 为等容过程，温度升高压强增大，$p_{c}>p_{a}$，故 $p_{b}<p_{c}$。两过程初、末温度相同，内能变化相同；过程 ab 中气体体积增大对外做功，过程 ac 中气体不做功，由热力学第一定律可知 $Q_{ab}>Q_{ac}$，C 正确。

（2）（i）考虑氢气的等温过程。该过程的初态压强为 $p_{0}$，体积为 $hS$，末态体积为 $0.8hS$。设末态的压强为 $p$，由玻意耳定律得 $p=\frac{p_{0}hS}{0.8hS}=1.25p_{0}$ ①。

活塞 A 从最高点被推回第一次平衡时位置的过程是等温过程。该过程的初态压强为 $1.1p_{0}$，体积为 $V$；末态的压强为 $p'$，体积为 $V'$，则 $p'=p+0.1p_{0}=1.35p_{0}$ ②，$V'=2.2hS$ ③。

由玻意耳定律得 $V=\frac{1.35p_{0}}{1.1p_{0}}\times2.2hS=2.7hS$ ④。

（ii）活塞 A 从最初位置升到最高点的过程为等压过程。该过程的初态体积和温度分别为 $2hS$ 和 $T_{0}=273\UK$，末态体积为 $2.7hS$。设末态温度为 $T$，由盖-吕萨克定律得 $T=\frac{2.7hS}{2hS}T_{0}=368.55\UK$ ⑤。
}

\item
%% number: 35
%% paperName: 2009年宁夏理综第 35 题 10 分
%% typeId: 6
%% type: 计算题
%% chapter: 光学
%% point: 光的折射
%% method: 分情况讨论
%% score: 10
%% degree: 450
%% duplicateId: 0
%% body:
\textbf{（1）}（5 分）某振动系统的固有频率为 $f_{0}$，在周期性驱动力的作用下做受迫振动，驱动力的频率为 $f$。若驱动力的振幅保持不变，下列说法正确的是\xzanswer{BD}（填入选项前的字母，有填错的不得分）。

\fourchoices[answer=BD]
{当 $f<f_{0}$ 时，该振动系统的振幅随 $f$ 增大而减小}
{当 $f>f_{0}$ 时，该振动系统的振幅随 $f$ 减小而增大}
{该振动系统的振动稳定后，振动的频率等于 $f_{0}$}
{该振动系统的振动稳定后，振动的频率等于 $f$}

\textbf{（2）}（10 分）一棱镜的截面为直角三角形 ABC，$\angle A=30^{\circ}$，斜边 AB$=a$。棱镜材料的折射率为 $n=\sqrt{2}$。在此截面所在的平面内，一条光线以 $45^{\circ}$ 的入射角从 AC 边的中点 M 射入棱镜射出的点的位置（不考虑光线沿原来路返回的情况）。

\onepicture[w=6.0cm, align=right]{figs/35.png}

%% answer:
\jdanswer{
\textbf{（1）}BD

\textbf{（2）}如果入射光线在法线的右侧，即出射点在 AB 边上离 A 点 $\frac{3}{8}a$ 的位置；如果入射光线在法线的左侧，出射点在 BC 边上离 B 点 $\frac{1}{8}a$ 的位置。
}
%% memo:
\memoanswer{
（1）该振动系统做受迫振动，振动稳定后，振动的频率等于驱动力的频率 $f$，与固有频率 $f_{0}$ 无关，故 D 正确，C 错误；由共振曲线可知，当 $f<f_{0}$ 时振幅随 $f$ 的增大而增大，当 $f>f_{0}$ 时振幅随 $f$ 的减小而增大，故 A 错误，B 正确。

（2）设入射角为 $i$，折射角为 $r$，由折射定律得 $\frac{\sin i}{\sin r}=n$ ①。

由已知条件及①式得 $r=30^{\circ}$ ②。

如果入射光线在法线的右侧，光路图如图 1 所示。设出射点为 F，由几何关系可得 $AF=\frac{3}{8}a$ ③。

即出射点在 AB 边上离 A 点 $\frac{3}{8}a$ 的位置。

\onepicture[w=6.0cm]{figs/35_an1.png}

如果入射光线在法线的左侧，光路图如图 2 所示。设折射光线与 AB 的交点为 D。

由几何关系可知，在 D 点的入射角 $\theta=60^{\circ}$ ④。

设全反射的临界角为 $\theta_{c}$，则 $\sin\theta_{c}=\frac{1}{n}$ ⑤。

由⑤和已知条件得 $\theta_{c}=45^{\circ}$ ⑥。

因此，光在 D 点全反射。

设此光线的出射点为 E，由几何关系得 $\angle DEB=90^{\circ}$，$BD=a-2AF$ ⑦，$BE=BD\sin30^{\circ}$ ⑧。

联立③⑦⑧式得 $BE=\frac{1}{8}a$ ⑨。

即出射点在 BC 边上离 B 点 $\frac{1}{8}a$ 的位置。

\onepicture[w=6.0cm]{figs/35_an2.png}
}

\item
%% number: 36
%% paperName: 2009年宁夏理综第 36 题 10 分
%% typeId: 6
%% type: 计算题
%% chapter: 运动和力
%% point: 动量和能量
%% method: 无
%% score: 10
%% degree: 450
%% duplicateId: 0
%% body:
\textbf{（1）}（5 分）关于光电效应，下列说法正确的是\xzanswer{A}（填入选项前的字母，有填错的不得分）。

\fourchoices[answer=A]
{极限频率越大的金属材料逸出功越大}
{只要光照射的时间足够长，任何金属都能产生光电效应}
{从金属表面出来的光电子的最大初动能越大，这种金属的逸出功越小}
{入射光的光强一定时，频率越高，单位时间内逸出的光电子数就越多}

\textbf{（2）}（10 分）两质量分别为 $M_{1}$ 和 $M_{2}$ 的劈 A 和 B，高度相同，放在光滑水平面上，A 和 B 的倾斜面都是光滑曲面，曲面下端与水平面相切，如图所示，一质量为 $m$ 的物块位于劈 A 的倾斜面上，距水平面的高度为 $h$。物块从静止滑下，然后又滑上劈 B。求物块在 B 上能够达到的最大高度。

\onepicture[align=right]{figs/36.png}

%% answer:
\jdanswer{
\textbf{（1）}A

\textbf{（2）}$h'=\frac{M_{1}M_{2}}{(M_{1}+m)(M_{2}+m)}h$
}
%% memo:
\memoanswer{
（1）由光电效应方程 $E_{k}=h\nu-W_{0}$ 及极限频率 $\nu_{c}=\frac{W_{0}}{h}$ 可知，极限频率越大的金属材料逸出功越大，A 正确；光电效应具有瞬时性，与照射时间无关，B 错误；光电子的最大初动能还与入射光的频率有关，不能仅由它判断逸出功大小，C 错误；入射光的光强一定时，频率越高，单位时间内逸出的光电子数不一定越多，D 错误。

（2）设物块到达劈 A 的低端时，物块和 A 的速度大小分别为 $v$ 和 $V$，由机械能守恒和动量守恒得 $mgh=\frac{1}{2}mv^{2}+\frac{1}{2}M_{1}V^{2}$ ①，$M_{1}V=mv$ ②。

设物块在劈 B 上达到的最大高度为 $h'$，此时物块和 B 的共同速度大小为 $V'$，由机械能守恒和动量守恒得 $mgh'+\frac{1}{2}(M_{2}+m)V'^{2}=\frac{1}{2}mv^{2}$ ③，$mv=(M_{2}+m)V'$ ④。

联立①②③④式得 $h'=\frac{M_{1}M_{2}}{(M_{1}+m)(M_{2}+m)}h$ ⑤。
}

\end{enumerate}

\end{document}
